% !TEX root = ../main.tex
\chapter{Introduction}\label{ch:introduction}

\section{Motivation and Goals}\label{sec:intro-motivation}

A weather forecast is useful only where it is applied, so the numerical fields that forecast or analyse weather have to capture global patterns while resolving local occurrences. Global reanalyses such as ERA5~\cite{hersbach2020era5} describe the atmosphere on a grid of roughly twenty-five kilometres, on which one whole valley with a mountain ridge in the alps is represented as one single number. No relevant local weather is decided at that scale. 

Closing that gap is called downscaling, and there are two established routes. The dynamical one runs a limited-area physical model at high resolution, taking its boundary conditions from the coarse field and letting the equations of motion generate the fine detail, as kilometre-scale models such as MeteoSwiss ICON-CH1 do~\cite{meteoswiss_icon}. It resolves Alpine structure far better than any reanalysis can, and is the operational standard in Switzerland, but running the ensemble it needs for confidence is computationally very expensive. The statistical approach learns the mapping from a period in which a coarse field and a fine observation both exist, then uses the mapping to infer local phenomena based on coarse fields alone. It buys the resolution for a fraction of the price and pays in a different currency: the mapping is only as good as the data it was fitted to, and carries no guarantee that its output is physically consistent. 

In this work we study the statistical approach to downscaling using machine learning approaches. That choice makes the Alps a demanding test domain. The terrain governing the local weather is the terrain the coarse model has averaged away, so the errors are not random noise to be filtered out but a systematic consequence of resolution, and the mapping has to reconstruct them from a field in which their cause is no longer visible. Switzerland is nevertheless unusually well placed for the attempt: MeteoSwiss maintains gridded daily analyses of temperature and precipitation on a one-kilometre lattice, interpolated from a dense station network~\cite{frei2014temperature, frei1998precip}, so there is a fine-scale reference that provides decades of "ground truth" data.


\section{Machine Learning for Downscaling: State of the Art}\label{sec:intro-sota}

Machine-learning approaches to statistical downscaling can be grouped along three axes: the input domain, the geometric mapping and the output.

On input, the field splits between enhancing the target variable, where the model refines a coarse version of the quantity it is predicting, and driving the prediction from the atmospheric state, where the surface variable is inferred from the broader meteorological dynamics above it (\sref{sec:bg-atmos-vs-surface}). A high-resolution elevation model, together with land-cover and seasonal encodings, is often added in either case to stand in for sub-grid processes. Training data comes almost universally from reanalysis~\cite{rampal2024enhancing}, in practice ERA5 or ERA5-Land~\cite{hersbach2020era5, munozsabater2021era5land}.

Geometrically there are two classes. Grid-to-grid architectures borrowed from image super-resolution~\cite{vandal2017deepsd} learn a mapping between two fixed lattices and have to be retrained when either of them changes. Point-based architectures instead model a continuous field that can be queried at any coordinate. \textcite{vaughan2022convcnp} takes the second route with convolutional conditional neural processes. 


Neural Processes (\sref{sec:bg-neural-processes}) have been used for off-grid environmental prediction and station-level downscaling~\cite{andersson2023convgnp}. On Swiss terrain specifically, \textcite{passos} adapted the approach of \textcite{vaughan2022convcnp} into a surface-driven ConvCNP variant compatible with the MeteoSwiss grid. That adaptation is the direct starting point of this work.

On output, the deficiencies of downscaling are well known and largely shared between approaches. Networks trained on a squared error regress towards the mean, which smooths fields spatially and clips the values for extreme observations; for precipitation specifically too many predicted rain days with little rain (drizzle) and missing the accurate prediction of heavy rain events are common. Generative approaches~\cite{harris2022generative} and parametric probabilistic heads are the two standard answers for these problems.

\textcite{cannon2008berngamma} follow the parametric route with a Bernoulli--Gamma density network. They report the improvement of heavy rain prediction by using a likelihood loss rather than a squared error. \textcite{mardani2025corrdiff} follow the generative route with a diffusion model that reaches kilometre resolution. They report the model can accurately reproduce the observed power spectrum, but that calibrating the predictive uncertainty remains unsolved. \textcite{bruinsma2023arcnp} draw joint samples from an already-trained conditional neural process by running it autoregressively which beats the pointwise performance of that same model. \textcite{rampal2022interpretable} found that when using synoptic drivers as input, the resulting model is relying on the meteorology rather than on geography, which is a significant step towards better interpretability. 

What none of these studies establishes is which atmospheric information a continuous-field model actually needs in complex terrain, or whether the reliance it develops is meteorological rather than mainly spatial and topographical. \textcite{rampal2022interpretable} come closest, showing for rainfall over the New Zealand Southern Alps that a grid-to-grid network driven by synoptic predictors relies on meteorology rather than on geography. Whether the same holds for a continuous-field model, and which atmospheric information such a model needs, has to our knowledge not been established for the European Alps.

\section{Research Questions}\label{sec:intro-questions}

\begin{itemize}
    \item \textbf{Feature selection and interpretability.} Can the accuracy and calibration of single-variable downscaling models be improved by using meteorologically relevant atmospheric predictors rather than surface-level data? And does the model learn physically meaningful relationships, or does it mainly act as a spatial interpolator? 
    \item \textbf{Computational efficiency.} How is computation cost coupled to input and architecture selection?
    What are the levers to reduce cost  without losing accuracy and robustness?
    \item \textbf{Temperature and Precipitation in complex terrain.} How well can the ConvCNP architecture be extended to downscale max temperature and precipitation accumulation over Switzerland, is it better than the results \textcite{vaughan2022convcnp} achieved on a European scale?
    \item \textbf{Mitigating the known weaknesses.} How far can drizzle bias, spatial oversmoothing and the underestimation of extremes be addressed by the choice of objective function and output distribution?
\end{itemize}
